\documentclass[11pt]{exam}
\usepackage{listings}
\usepackage{pdfsync}
\usepackage{amsmath}
\usepackage{color}


%
%  Created by Nancy Ide on 2008-12-10.
%  Copyright (c) 2008 __MyCompanyName__. All rights reserved.
%
%

% This is now the recommended way for checking for PDFLaTeX:
\usepackage{ifpdf}


%\newif\ifpdf
%\ifx\pdfoutput\undefined
%\pdffalse % we are not running PDFLaTeX
%\else
%\pdfoutput=1 % we are running PDFLaTeX
%\pdftrue
%\fi

\ifpdf
\usepackage{subfigure}
\usepackage[pdftex]{graphicx}
\else
\usepackage{graphicx}
\fi
\usepackage{enumerate}

%
%  Update these values for running headers
%

%\runningheader{\bf Departure Checklist}{}{\bf Revised 2015-03-26}

\begin{document}
\begin{center}

{\bf\Large CHCCA Opening and Closing Procedures}\\
\textit{10/10/21}
\end{center}
\vskip.25in

\section*{Closing}
 
\begin{itemize}
\item Contact the Property Manager (PM) in July or early August to inform him of our desired close date (usually within the first week of November).
The closing date is highly dependent on the plumber availability.  
\item Ensure that the PM and plumber understand that the close date for cottages 18 and19 is usually later in the month.  
\item Obtain desired close date from 18 and19 and inform PM. 
\item Ensure that PM coordinates with Days for propane shutoff.
\item Communicate close date to owners via email and at the owners meeting. 
\item  At our annual meeting ask owners to communicate to the PM any specific close requests for their individual cottages and the date they will be leaving their cottages for the season by 10/31. This allows the PM to get a head start on individual cottage close activities.
\end{itemize}  
 
\subsection*{Property Manager Closing Activities}  
 
\begin{itemize}
\item Clean and store trash barrels
\item Store BBQ grills
\item Remove private beach signs from beach and place in storage room
\item Remove and dispose of any beach paraphernalia left on beach (remove/dispose in November) (secured Kayaks/boats can stay on the beach)
\item Shutter all cottage doors and windows except for cottage10 SLW
\item Store Adirondack chairs in each individual cottage.
\item Unplug stove, refrigerator, lights, heaters and all small appliances within each cottage
\item Board up laundry shed
\item Leave electricity on for all cottages
\item Shut electricity off for laundry shed
\item Adjust exterior light timers  (lights on 4PM and off at 6AM)
\item Put garden hoses in the storage room 
\end{itemize}  
 
\subsection*{Plumber activities: (PM to coordinate)}
\begin{itemize}
\item Drain and blow and winterize CHCCA water pipes and turn off well pumps in the well pit.
\item Drain, blow and winterize all individual cottage water pipes including sinks, toilet and stove, except for cottage 10 SLW.
\item Drain, blow and winterize the laundry shed.
 \end{itemize}  

\subsection*{Days Propane (PM to coordinate)}
\begin{itemize}
\item Notify Days of CHCCA close date.
 \end{itemize}  
 
\subsection*{Individual cottage close procedures}
\begin{itemize}
\item Each cottage should inform PM of any additional close activities for their cottage ( e.g., put mouse traps down, remove pictures from wall, remove screens etc). PM to invoice owner as needed.
  \end{itemize}  
  
\subsection*{Seaside Disposal (Brice to coordinate)}
\begin{itemize}
\item Inform Scott Elliott of close date via email (ssdisposal@verizon.net)so that he can collect recycle bins and stop trash collection.

   \end{itemize}  
 
\section*{Opening}
 
\begin{itemize}
\item Immediately after closing, work with the PM to secure an opening date (usually within the first week of April).  Note that the opening date is different than the occupancy date: the opening date is the date that the water is turned on so that it can be tested; the occupancy date is established once we receive the water testing results.
 
\item As soon as the opening date is established, contact WhiteWater, Inc.  (whitewateronline.com) so that they can be on site to take the water samples.  Start by contacting Alisha Almeida no later than March 1 (earlier is better) to schedule a testing date.

WhiteWater contact info:
\begin{itemize}
\item Alisha Almeida, scheduling coordinator,- 508-248-2892/AAlmeida@rhwhite.com
\item Jonathan Preissler, operations supervisior, 508-612-5150/JPreissler@rhwhite.com
\item John, water tester, 508-864-0352 
   \end{itemize} 
    
\item \textbf{Water testing}: The well pit and at least two cottages need to have running water in order for WhiteWater to do the water sampling.  So, for example, if the plumber is turning the water on the morning of April 3rd, try to schedule the water testing that afternoon or the following day.  The PM needs to be onsite for the water testing to be sure the water tester has access to the well pit. It usually take about five days to get the test results back and the DEP certificate. 

 
\item Once the water test results are received, forward to the webmaster for posting on our website and communicate the official occupancy date to the owners, stressing that 
the DEP is very clear that occupancy cannot occur to until  the water testing results are received; otherwise, 
the Association may be fined, and the individual owners who violated the dictate will be asked to cover the cost.
 
\item The rest of the opening activities are the same as closing except in reverse.
   \end{itemize}
    
\end{document}
